\section{The contrapositive and proof by contradiction}\label{ch04:sec:contrapositive}
Now turn the claim of Section~\ref{ch04:sec:directions} around. Its converse (Section~\ref{ch02:sec:implication}) asks: if $n^2$ is even, must the integer $n$ be even? Working forward gives $n^2=2w$ for some integer $w$, which describes $n^2$, not $n$; taking square roots gives only $n=\sqrt{2w}$ or $n=-\sqrt{2w}$, neither visibly twice an integer. The forward direction stalls.

So try the contrapositive (Section~\ref{ch02:sec:contrapositive}):
\begin{center}
\emph{If $n$ is not even, then $n^2$ is not even.}
\end{center}
Since an implication and its contrapositive are equivalent, proving this proves the converse. But the new hypothesis, ``$n$ is not even,'' says only what $n$ is \emph{not}. If we knew that such an $n$ must be odd, we could write $n=2k+1$ and square it. Section~\ref{ch01:sec:definition} stated without proof that every integer is either even or odd, but not both. We now prove this, using a method called \emph{proof by contradiction}.

\subsection*{Proof by contradiction}
To prove a claim by contradiction, temporarily suppose that the claim is false, keep all the original hypotheses, and reason until you reach a \emph{contradiction}: some statement and its negation both follow. If the supposition were true, everything correctly deduced from it and the hypotheses would be true as well; but a statement and its negation cannot both be true. So the supposition is false. This last step uses the principle that every statement is either true or false (Section~\ref{ch02:sec:contrapositive}): ruling out ``the claim is false'' leaves only ``the claim is true.''

Section~\ref{ch02:sec:contrapositive} already used this method, in its discussion of vacuous truth, to show that \emph{no integer is both even and odd}. If some integer $n$ were both, then $n=2k$ and $n=2\ell+1$ for integers $k$ and $\ell$, so $2(k-\ell)=1$, and the integer $k-\ell$ would equal $1/2$, which is not an integer. A statement (``$k-\ell$ is an integer'') and its negation would both hold, so the supposition is impossible.

The end point must be a genuine impossibility, not something that only ``looks unlikely.'' Mark the supposition clearly, for example with ``suppose, for contradiction,'' because it is temporary: inside the argument we may use it freely, but once the contradiction is reached it is discarded, and nothing deduced from it counts as established afterward.

\subsection*{Every integer is even or odd}
The other half, that every integer is even or odd, needs one more fact, which we accept without proof; it is the third of the starting facts listed in \emph{To the reader}.
\begin{center}
\emph{Well-ordering property: every nonempty set of positive integers has a smallest member.}
\end{center}
For example, the positive integers whose square is larger than $50$ form a nonempty set, since it contains $10$. Its smallest member is $8$: the squares of $1,2,\ldots,7$ are at most $49$, while $8^2=64$. Each part of the property matters. The empty set has no smallest member, since it has no members. The negative integers have none either, since each negative integer $n$ has a smaller negative integer, $n-1$; nor do the positive rational numbers, since each positive rational $r$ has a smaller positive rational, $r/2$.

\par\noindent\begin{minipage}{\linewidth}
\begin{lemma}\label{ch04:lem:even-or-odd}
Every integer is either even or odd, and no integer is both.
\end{lemma}
\end{minipage}\par

\begin{proof}
We showed above that no integer is both even and odd. It remains to show that every integer $n$ is even or odd. If $n=0$, then $n$ is even, since $0=2\cdot0$.

Next, suppose, for contradiction, that some positive integer is neither even nor odd. Then the set of such positive integers is nonempty, so by the well-ordering property it has a smallest member; call it $m$. Now $m\ne1$, because $1=2\cdot0+1$ is odd. So $m\ge2$, and $m-1$ is a positive integer smaller than $m$. Since $m$ is the \emph{smallest} positive integer that is neither even nor odd, $m-1$ is even or odd. If $m-1=2k$ for an integer $k$, then $m=2k+1$, so $m$ is odd. If $m-1=2k+1$ for an integer $k$, then $m=2k+2=2(k+1)$, so $m$ is even. Either way $m$ is even or odd, which contradicts the choice of $m$. Hence every positive integer is even or odd.

Finally, let $n$ be negative. Then $-n$ is positive, so by the previous paragraph $-n=2k$ or $-n=2k+1$ for some integer $k$. In the first case $n=2(-k)$ is even. In the second, $n=-2k-1=2(-k-1)+1$ is odd; for instance, $7=2\cdot3+1$ and $-7=2\cdot(-4)+1$.
\end{proof}

Here the contradiction is that $m$ has a property it was chosen \emph{not} to have. Proofs that begin ``choose the smallest \ldots'' often end this way: the next section chooses a smallest denominator, and Chapter~\ref{ch:induction} uses the same pattern to explain why induction works.

\subsection*{Back to the contrapositive}
Lemma~\ref{ch04:lem:even-or-odd} lets us finish the argument. Since the next section uses the result twice, we record it as a lemma.

\par\noindent\begin{minipage}{\linewidth}
\begin{lemma}\label{thm:parity-square}
For every integer $n$, the integer $n$ is even if and only if $n^2$ is even.
\end{lemma}
\end{minipage}\par

\begin{proof}
Let $n$ be an integer. If $n$ is even, then $n^2$ is even, by the direct proof in Section~\ref{ch04:sec:directions}.

For the other direction we prove the contrapositive: if $n$ is not even, then $n^2$ is not even. So suppose $n$ is not even. By Lemma~\ref{ch04:lem:even-or-odd}, $n$ is odd, so $n=2k+1$ for some integer $k$. Then
\begin{align*}
n^2&=(2k+1)(2k+1)\\
&=4k^2+4k+1\\
&=2(2k^2+2k)+1.
\end{align*}
Since $2k^2+2k$ is an integer, $n^2$ is odd, so by Lemma~\ref{ch04:lem:even-or-odd} again, $n^2$ is not even. This proves the contrapositive, and hence that if $n^2$ is even, then $n$ is even.
\end{proof}

In general, try the contrapositive when the negation of the conclusion is easier to unpack than the hypothesis; when the hypothesis is the easier one, as in Section~\ref{ch04:sec:directions}, prove the statement directly.

\pauseq{The proof of Lemma~\ref{thm:parity-square} used Lemma~\ref{ch04:lem:even-or-odd} twice. Which half of Lemma~\ref{ch04:lem:even-or-odd} was used each time?}{The first use needs ``every integer is even or odd'': it turns ``$n$ is not even'' into ``$n$ is odd,'' which gives $n=2k+1$. The second needs ``no integer is both'': it turns ``$n^2$ is odd'' into ``$n^2$ is not even.''}
